\chapter{Hochschild homology}\label{ch:hochschild}
\providecommand{\cA}{\mathcal{A}}
\providecommand{\cB}{\mathcal{B}}

The algebra $\Psido{\Z}(X)$ of classical pseudodifferential operators of
integral order on a compact manifold $X$, studied in \cref{ch:manifolds}, is
a non-commutative deformation of a commutative algebra: modulo
$\Psido{-\infty}(X)$ it is, via the full symbol, the algebra of formal sums of
homogeneous functions on $T^*X\setminus0$, with a product whose leading term
is pointwise multiplication and whose first correction is the Poisson bracket.
Hochschild homology is the basic homological invariant of an associative
algebra; in degree zero it is the dual of the space of traces. The aim of this
chapter is to compute the Hochschild homology of the algebra
$\Psido{\Z}(X)/\Psido{-\infty}(X)$ of formal symbols, following
Brylinski--Getzler \cite{Brylinski-Getzler} and Wodzicki
\cite{Wodzicki-residue,Wodzicki-cyclic}:
\begin{equation*}
  \HH_k\bigl(\Psido{\Z}(X)/\Psido{-\infty}(X)\bigr)\cong H^{2n-k}(S^*X)\oplus
  H^{2n-k-1}(S^*X),\qquad n=\dim X .
\end{equation*}
The computation proceeds in three stages. First (\cref{hochschild:sec:complex,hochschild:sec:polynomial})
we set up the Hochschild complex and compute it for polynomial algebras by
means of the Koszul resolution. Secondly (\cref{hochschild:sec:smooth,hochschild:sec:commutative})
we prove the theorem of Hochschild--Kostant--Rosenberg and Connes that the
Hochschild homology of $\CI(X)$ is the space of differential forms, and deduce
the analogous statement for the commutative algebra of homogeneous symbols.
Finally (\cref{hochschild:sec:symbols,hochschild:sec:spectral,hochschild:sec:E2,hochschild:sec:main})
we filter the Hochschild complex of the symbol algebra by order; the
resulting spectral sequence has $E_1$ term the homogeneous differential forms
on $T^*X\setminus0$, its first differential is the Poisson (Koszul--Brylinski)
differential, a symplectic Hodge star identifies $E_2$ with de Rham cohomology,
and a homogeneity argument shows that the spectral sequence degenerates and
converges.

Throughout, vector spaces are complex, $\otimes$ without decoration is the
algebraic tensor product over $\C$, and $X$ is a smooth manifold of dimension
$n$.

\section{The Hochschild complex}\label{hochschild:sec:complex}

\subsection{Chains and differentials}
Let $\cA$ be an associative algebra over $\C$, not necessarily with identity.
The space of (algebraic) Hochschild $k$-chains, $k\in\N_0$, is
\begin{equation*}
  C_k(\cA)=\cA^{\otimes(k+1)}=\cA\otimes\cA\otimes\dots\otimes\cA .
\end{equation*}
For $k\ge1$ define the \emph{face maps} $d_i:C_k(\cA)\longrightarrow C_{k-1}(\cA)$,
$0\le i\le k$, by
\begin{equation}
\begin{aligned}
  d_i(a_0\otimes\dots\otimes a_k)&=a_0\otimes\dots\otimes a_{i-1}\otimes
  a_ia_{i+1}\otimes a_{i+2}\otimes\dots\otimes a_k,\quad 0\le i\le k-1,\\
  d_k(a_0\otimes\dots\otimes a_k)&=a_ka_0\otimes a_1\otimes\dots\otimes a_{k-1},
\end{aligned}
\label{hochschild:eq:faces}
\end{equation}
and set
\begin{equation}
  b'=\sum_{i=0}^{k-1}(-1)^id_i,\qquad b=b'+(-1)^kd_k=\sum_{i=0}^k(-1)^id_i
  \quad\text{on }C_k(\cA),
  \label{hochschild:eq:b}
\end{equation}
with $b=b'=0$ on $C_0(\cA)=\cA$.

\begin{lemma}\label{hochschild:lem:simplicial}
For $k\ge2$ and $0\le i<j\le k$ the face maps satisfy
$d_id_j=d_{j-1}d_i$ as maps $C_k(\cA)\longrightarrow C_{k-2}(\cA)$.
Consequently $(b')^2=0$ and $b^2=0$.
\end{lemma}
\begin{proof}
If $j\le k-1$ both sides only multiply neighbouring factors and the identity
is the associativity of the product (for $j=i+1$ both sides replace
$a_i\otimes a_{i+1}\otimes a_{i+2}$ by $a_ia_{i+1}a_{i+2}$; for $j>i+1$ the
two multiplications commute). If $j=k$ and $i=0$, both sides give
$a_ka_0a_1\otimes a_2\otimes\dots\otimes a_{k-1}$; if $j=k$ and $0<i<k-1$,
both give $a_ka_0\otimes\dots\otimes a_ia_{i+1}\otimes\dots\otimes a_{k-1}$;
and if $j=k$, $i=k-1$, both give $a_{k-1}a_ka_0\otimes a_1\otimes\dots\otimes
a_{k-2}$.

Now $b^2=\sum_{i=0}^{k-1}\sum_{j=0}^{k}(-1)^{i+j}d_id_j$ on $C_k(\cA)$. The
map $(i,j)\longmapsto(j-1,i)$ is a bijection from the pairs with $i<j\le k$
onto the pairs $(i',j')$ with $j'\le i'\le k-1$, and by the identity the
corresponding terms are equal but carry the opposite signs $(-1)^{i+j}$ and
$(-1)^{i+j-1}$; hence $b^2=0$. The same argument with $k$ replaced by $k-1$
in the range of $j$ gives $(b')^2=0$.
\end{proof}

\begin{definition}\label{hochschild:def:HH}
The \emph{Hochschild homology} of $\cA$ is the homology of the complex
$(C_\bullet(\cA),b)$:
\begin{equation*}
  \HH_k(\cA)=\frac{\Nul\bigl(b:C_k(\cA)\to C_{k-1}(\cA)\bigr)}
  {\Ran\bigl(b:C_{k+1}(\cA)\to C_k(\cA)\bigr)},\qquad k\in\N_0 .
\end{equation*}
\end{definition}

In degree zero $b(a_0\otimes a_1)=a_0a_1-a_1a_0$, so
$\HH_0(\cA)=\cA/[\cA,\cA]$, where $[\cA,\cA]$ is the span of the commutators.
Thus the linear functionals on $\HH_0(\cA)$ are precisely the \emph{traces}
on $\cA$, the linear maps $\tau$ with $\tau(ab)=\tau(ba)$. If $\cA$ is
commutative then $b=0$ on $C_1(\cA)$ and $\HH_0(\cA)=\cA$.

\subsection{The bar complex}
The complex $(C_\bullet(\cA),b')$ is the \emph{bar complex}. It is, in
contrast to the Hochschild complex, trivial for algebras with identity.

\begin{lemma}\label{hochschild:lem:bar-acyclic}
If $\cA$ has an identity $1$, then $s(a_0\otimes\dots\otimes a_k)=1\otimes
a_0\otimes\dots\otimes a_k$ defines maps $s:C_k(\cA)\to C_{k+1}(\cA)$ with
\begin{equation*}
  b's+sb'=\Id\quad\text{on }C_k(\cA),\ k\ge0;
\end{equation*}
in particular $(C_\bullet(\cA),b')$ is acyclic.
\end{lemma}
\begin{proof}
On $C_{k+1}(\cA)$, $d_0(1\otimes a)=a$ for $a\in C_k(\cA)$, while for $1\le
i\le k$, $d_i(1\otimes a)=1\otimes d_{i-1}a$. Hence
$b'sa=a+\sum_{i=1}^{k}(-1)^i\,1\otimes d_{i-1}a=a-sb'a$.
\end{proof}

\begin{definition}\label{hochschild:def:H-unital}
An associative algebra is \emph{H-unital} if its bar complex
$(C_\bullet(\cA),b')$ is acyclic.
\end{definition}

So every unital algebra is H-unital. For H-unital algebras Hochschild homology
has good excision properties (Wodzicki); we shall not need these.

\subsection{Hochschild homology as a derived functor}
\label{hochschild:sec:Tor}
Suppose $\cA$ has an identity and let $\cA^e=\cA\otimes\cA^{\mathrm{op}}$, so that
$\cA$-bimodules are the same as left $\cA^e$-modules; $\cA$ itself is an
$\cA^e$-module by $(a\otimes a')\cdot c=aca'$. Set
\begin{equation*}
  P_k=\cA\otimes\cA^{\otimes k}\otimes\cA\ (=\cA^{\otimes(k+2)}),
  \qquad (a\otimes a')\cdot(a_0\otimes\dots\otimes a_{k+1})
  =aa_0\otimes a_1\otimes\dots\otimes a_{k+1}a'.
\end{equation*}
Each $P_k\cong\cA^e\otimes\cA^{\otimes k}$ is a free $\cA^e$-module, the bar
differential $b'$ of \eqref{hochschild:eq:b} (on $\cA^{\otimes(k+2)}$, i.e.\
$b'=\sum_{i=0}^{k}(-1)^id_i$) is $\cA^e$-linear, and by
\cref{hochschild:lem:bar-acyclic} the complex
\begin{equation}
  \dots\xrightarrow{b'}P_2\xrightarrow{b'}P_1\xrightarrow{b'}P_0\xrightarrow{\ m\ }\cA\longrightarrow0,
  \qquad m(a_0\otimes a_1)=a_0a_1,
  \label{hochschild:eq:bar-res}
\end{equation}
is exact: it is a free resolution of $\cA$ as an $\cA^e$-module. Applying
$\cA\otimes_{\cA^e}(\cdot)$ and using the isomorphism
\begin{equation*}
  \cA\otimes_{\cA^e}P_k\longrightarrow C_k(\cA),\quad
  c\otimes(a_0\otimes a_1\otimes\dots\otimes a_k\otimes a_{k+1})\longmapsto
  a_{k+1}ca_0\otimes a_1\otimes\dots\otimes a_k,
\end{equation*}
the differential $\Id\otimes b'$ becomes $b$: the term $d_0$ of $b'$ gives
$d_0$, and the last term $(-1)^k(\dots\otimes a_ka_{k+1})$ of $b'$ on $P_k$
(with $a_{k+1}=1$) gives $(-1)^ka_kca_0\otimes\dots=(-1)^kd_k$. Thus
\begin{equation*}
  \HH_k(\cA)=\operatorname{Tor}^{\cA^e}_k(\cA,\cA).
\end{equation*}
The practical content of this is the following standard fact
(the \emph{comparison theorem} of homological algebra): if
$Q_\bullet\to\cA\to0$ is another resolution of $\cA$ by free (or projective)
$\cA^e$-modules, then there are $\cA^e$-linear chain maps $\Phi:Q_\bullet\to
P_\bullet$ and $\Psi:P_\bullet\to Q_\bullet$ inducing the identity on $\cA$,
and $\Psi\Phi$, $\Phi\Psi$ are $\cA^e$-linearly chain homotopic to the
identity. The maps are constructed degree by degree: on a basis element $e$
of $Q_k$ one chooses $\Phi(e)$ in $P_k$ with $b'\Phi(e)=\Phi(\partial e)$,
which is possible because $\Phi(\partial e)$ is a cycle (for $k=0$, lies over
the right element of $\cA$) and $P_\bullet$ is exact; the homotopies are
constructed in the same way. Consequently $\HH_\bullet(\cA)$ is also the
homology of $\cA\otimes_{\cA^e}Q_\bullet$.

\subsection{Topological algebras}\label{hochschild:sec:topological}
For the `large' algebras of analysis the algebraic tensor product is
inappropriate and is replaced by a completion. If $\cA$ is a Fréchet algebra
(multiplication jointly continuous) one uses the completed projective tensor
product $\cA^{\hat\otimes(k+1)}$; the face maps extend by continuity and
\cref{hochschild:lem:simplicial,hochschild:lem:bar-acyclic} continue to hold,
by density of the algebraic tensor product. The resulting homology is the
\emph{continuous} Hochschild homology.

For $\cA=\CI(X)$, $X$ compact, a theorem of Grothendieck identifies
$\CI(X)\hat\otimes\CI(Y)$ with $\CI(X\times Y)$ (this rests on the nuclearity
of $\CI(X)$; the algebraic tensor product is in any case dense in
$\CI(X\times Y)$). We shall simply \emph{define} the chain spaces for $\CI(X)$
to be $\CI(X^{k+1})$, see \cref{hochschild:sec:smooth}; for the symbol
algebras the completion used is specified in \cref{hochschild:def:chains-comm,hochschild:def:chains-full}.

\section{Polynomial algebras}\label{hochschild:sec:polynomial}

Let $\cA=\C[x]=\C[x_1,\dots,x_n]$, regarded as the algebra of polynomial
functions on $\R^n$. Then $C_k(\cA)=\C[z_0,z_1,\dots,z_k]$, $z_j\in\R^n$, is
the space of polynomials on $(\R^n)^{k+1}$, the element
$a_0\otimes\dots\otimes a_k$ corresponding to $a_0(z_0)a_1(z_1)\cdots
a_k(z_k)$. In these terms the face maps \eqref{hochschild:eq:faces} are
restrictions to diagonals:
\begin{equation}
\begin{aligned}
  (d_iu)(z_0,\dots,z_{k-1})&=u(z_0,\dots,z_{i-1},z_i,z_i,z_{i+1},\dots,z_{k-1}),
  \quad 0\le i\le k-1,\\
  (d_ku)(z_0,\dots,z_{k-1})&=u(z_0,z_1,\dots,z_{k-1},z_0).
\end{aligned}
\label{hochschild:eq:faces-fn}
\end{equation}
Let $\Omega^k_{\mathrm{pol}}=\C[x]\otimes\Lambda^k(\C^n)$ be the space of
$k$-forms on $\R^n$ with polynomial coefficients, $dx^j$ corresponding to the
$j$-th basis vector $e_j$ of $\C^n$.

\subsection{The Hochschild--Kostant--Rosenberg map}
For a function $u$ of $(z_0,\dots,z_k)$ set
\begin{equation}
  \varepsilon(u)(x)=\frac1{k!}\sum_{j_1,\dots,j_k=1}^n
  \frac{\pa^ku}{\pa z_1^{j_1}\cdots\pa z_k^{j_k}}(x,x,\dots,x)\,
  dx^{j_1}\wedge\dots\wedge dx^{j_k},
  \label{hochschild:eq:HKR}
\end{equation}
where $z_l^{j}$ is the $j$-th coordinate of $z_l$. Thus
$\varepsilon(a_0\otimes a_1\otimes\dots\otimes a_k)=\frac1{k!}a_0\,da_1\wedge\dots\wedge da_k$.
For $a_0,\dots,a_k\in\cA$ we also set
\begin{equation}
  \varphi(a_0;a_1,\dots,a_k)=\sum_{\sigma\in\mathfrak S_k}\sgn(\sigma)\,
  a_0\otimes a_{\sigma(1)}\otimes\dots\otimes a_{\sigma(k)}\in C_k(\cA).
  \label{hochschild:eq:antisym}
\end{equation}

\begin{lemma}\label{hochschild:lem:HKR-chain}
For any commutative algebra $\cA$, define
$\varepsilon:C_k(\cA)\to\Omega^k_{\cA}$ by $\varepsilon(a_0\otimes\dots\otimes
a_k)=\frac1{k!}a_0\,da_1\wedge\dots\wedge da_k$, where $\Omega^\bullet_\cA$
is any graded-commutative algebra containing $\cA$ in degree zero, with a
derivation $d$ (for instance the de Rham complex of a manifold, with $\cA$ a
subalgebra of the smooth functions). Then $\varepsilon\circ b=0$,
$b\,\varphi(a_0;a_1,\dots,a_k)=0$ and
$\varepsilon\bigl(\varphi(a_0;a_1,\dots,a_k)\bigr)=a_0\,da_1\wedge\dots\wedge da_k$.
\end{lemma}
\begin{proof}
Write $T_i=a_0a_i\,da_1\wedge\dots\wedge\widehat{da_i}\wedge\dots\wedge da_k$.
By the Leibniz rule $d(a_ia_{i+1})=a_{i+1}da_i+a_ida_{i+1}$, so
$(k-1)!\,\varepsilon(b(a_0\otimes\dots\otimes a_k))$ equals
\begin{equation*}
  T_1+\sum_{i=1}^{k-1}(-1)^i(T_{i+1}+T_i)+(-1)^kT_k=0 .
\end{equation*}
In $b\varphi$, the terms $d_i$ with $1\le i\le k-1$ cancel in pairs, since
$\sigma$ and $\sigma\circ(i\ i{+}1)$ have opposite signs and $\cA$ is
commutative. The term $d_0$ for $\sigma$ is $\sgn(\sigma)a_0a_{\sigma(1)}\otimes
a_{\sigma(2)}\otimes\dots$, and the term $(-1)^kd_k$ for $\sigma'$, where
$\sigma'(j)=\sigma(j+1)$ for $j<k$ and $\sigma'(k)=\sigma(1)$, is
$(-1)^k\sgn(\sigma')a_{\sigma(1)}a_0\otimes a_{\sigma(2)}\otimes\dots$; since
$\sgn(\sigma')=(-1)^{k-1}\sgn(\sigma)$ these cancel. Finally
$\varepsilon\varphi=\frac1{k!}\sum_\sigma\sgn(\sigma)a_0\,da_{\sigma(1)}\wedge\dots\wedge
da_{\sigma(k)}=a_0\,da_1\wedge\dots\wedge da_k$.
\end{proof}

\subsection{The Koszul resolution}
We identify $\cA^e=\C[x]\otimes\C[x]$ with $\C[x,y]$, polynomials in
$(x,y)\in\R^n\times\R^n$, via $a\otimes a'\mapsto a(x)a'(y)$; then $\cA$ is the
$\cA^e$-module $\C[x,y]/(x_1-y_1,\dots,x_n-y_n)$. The \emph{Koszul complex} is
\begin{equation*}
  K_k=\C[x,y]\otimes\Lambda^k(\C^n),\quad
  \pa(f\,e_{i_1}\wedge\dots\wedge e_{i_k})=\sum_{l=1}^k(-1)^{l-1}
  (x_{i_l}-y_{i_l})f\,e_{i_1}\wedge\dots\wedge\widehat{e_{i_l}}\wedge\dots\wedge e_{i_k},
\end{equation*}
with augmentation $K_0=\C[x,y]\to\cA$, $f\mapsto f(x,x)$.

\begin{lemma}\label{hochschild:lem:Koszul}
The augmented Koszul complex $\dots\to K_1\to K_0\to\cA\to0$ is exact, so $K_\bullet$
is a free resolution of $\cA$ as an $\cA^e$-module.
\end{lemma}
\begin{proof}
In the variables $u=x-y$, $v=y$ identify $K_k$ with the $k$-forms in $u$ with
coefficients in $\C[u,v]$, $e_i\leftrightarrow du_i$. Then $\pa=\iota_E$ is
contraction with the Euler field $E=\sum_iu_i\pa_{u_i}$ and the augmentation
is $f\mapsto f|_{u=0}$. Let $d_u$ be the exterior derivative in $u$ (with $v$
as a parameter). By Cartan's formula $d_u\iota_E+\iota_Ed_u=\mathcal L_E$,
and $\mathcal L_E$ multiplies $u^\alpha du_I$ (with coefficients in $\C[v]$)
by $|\alpha|+|I|$. Hence if $\omega\in K_k$, $k\ge1$, is homogeneous of total
degree $N=|\alpha|+|I|\ge1$ in $u$ and $\pa\omega=0$, then
$\omega=\pa(d_u\omega/N)$; splitting a general cycle into homogeneous parts,
$K_\bullet$ is exact in degrees $k\ge1$. In degree $0$ the kernel of the
augmentation consists of the $f$ with no term of degree $0$ in $u$, and such
an $f$ equals $\pa(\sum_N d_uf_N/N)$, $f_N$ its homogeneous parts. Finally the
augmentation is surjective.
\end{proof}

\begin{theorem}[Hochschild--Kostant--Rosenberg]\label{hochschild:thm:polynomial}
For the polynomial algebra $\cA=\C[x_1,\dots,x_n]$ the map
\eqref{hochschild:eq:HKR} induces an isomorphism
\begin{equation*}
  \varepsilon_*:\HH_k(\C[x])\longrightarrow\Omega^k_{\mathrm{pol}}
  =\C[x]\otimes\Lambda^k(\C^n),\qquad k\in\N_0,
\end{equation*}
with inverse induced by $a_0\,dx^{i_1}\wedge\dots\wedge dx^{i_k}\longmapsto
\varphi(a_0;x_{i_1},\dots,x_{i_k})$. In particular $\HH_k(\C[x])=0$ for
$k>n$.
\end{theorem}
\begin{proof}
Let $P_\bullet$ be the bar resolution \eqref{hochschild:eq:bar-res} and define
the $\cA^e$-linear maps $\Phi:K_k\to P_k$ by
\begin{equation*}
  \Phi(e_{i_1}\wedge\dots\wedge e_{i_k})=\sum_{\sigma\in\mathfrak S_k}\sgn(\sigma)\,
  1\otimes x_{i_{\sigma(1)}}\otimes\dots\otimes x_{i_{\sigma(k)}}\otimes1 .
\end{equation*}
This is a chain map. Indeed, applying $b'$, the inner terms (multiplying two
adjacent $x$'s) cancel in pairs as in \cref{hochschild:lem:HKR-chain}. The
first term is $x_{i_{\sigma(1)}}\cdot(1\otimes x_{i_{\sigma(2)}}\otimes\dots\otimes1)$
(left action, i.e.\ multiplication by the variable $x$) and, grouping the
$\sigma$ with $\sigma(1)=l$, contributes
$\sum_l(-1)^{l-1}x_{i_l}\Phi(e_{i_1}\wedge\dots\widehat{e_{i_l}}\dots\wedge e_{i_k})$;
the last term is $(-1)^k(1\otimes\dots\otimes1)\cdot x_{i_{\sigma(k)}}$
(right action, multiplication by $y$) and, grouping the $\sigma$ with
$\sigma(k)=l$ (sign $(-1)^{k-l}$), contributes
$-\sum_l(-1)^{l-1}y_{i_l}\Phi(\dots\widehat{e_{i_l}}\dots)$. Together this is
$\Phi(\pa e_{i_1}\wedge\dots\wedge e_{i_k})$. In degree $0$, $\Phi$ is the
identity of $\cA^e$, so $\Phi$ lies over the identity of $\cA$.

By \cref{hochschild:lem:Koszul} both complexes are free resolutions of $\cA$,
so by the comparison theorem (\cref{hochschild:sec:Tor}) $\Id\otimes\Phi$
induces an isomorphism from the homology of $\cA\otimes_{\cA^e}K_\bullet$ to
that of $\cA\otimes_{\cA^e}P_\bullet$, i.e.\ to $\HH_\bullet(\cA)$. Now
$\cA\otimes_{\cA^e}K_k=\cA\otimes\Lambda^k(\C^n)=\Omega^k_{\mathrm{pol}}$ and
the induced differential vanishes, since $x_i-y_i$ acts on $\cA$ as $0$. Under
the identification $\cA\otimes_{\cA^e}P_k=C_k(\cA)$, $\Id\otimes\Phi$ becomes
$a_0\,dx^{i_1}\wedge\dots\wedge dx^{i_k}\mapsto\varphi(a_0;x_{i_1},\dots,x_{i_k})$,
which therefore induces an isomorphism $\varphi_*:\Omega^k_{\mathrm{pol}}\to\HH_k(\cA)$.
By \cref{hochschild:lem:HKR-chain}, $\varepsilon$ descends to homology and
$\varepsilon_*\varphi_*=\Id$, so $\varepsilon_*=\varphi_*^{-1}$.
\end{proof}

\begin{remark}\label{hochschild:rem:Laurent}
The same argument applies to the algebra $\C[r,r^{-1}]$ of Laurent
polynomials in one variable. Here $\cA^e=\C[r^{\pm1},s^{\pm1}]$ and the Koszul
resolution is $0\to\cA^e\xrightarrow{\ r-s\ }\cA^e\to\cA\to0$; it is exact
because $r-s$ is not a zero divisor and $f(r,s)-f(r,r)$ is divisible by $r-s$
in $\cA^e$ (write $f=(rs)^{-N}g$ with $g$ a polynomial). Hence
$\HH_0(\C[r^{\pm1}])=\C[r^{\pm1}]$, $\HH_1(\C[r^{\pm1}])=\C[r^{\pm1}]\,dr$
and $\HH_k(\C[r^{\pm1}])=0$ for $k\ge2$, the isomorphisms again being
induced by $\varepsilon$.
\end{remark}

The original direct argument for \cref{hochschild:thm:polynomial} proceeds
instead by successively restricting to chains vanishing on diagonals. This
reduction is very useful for smooth functions, where no finite free
resolution is at hand, and we now describe it.

\section{Smooth functions}\label{hochschild:sec:smooth}

\subsection{The complex}
Let $X$ be a compact manifold. We take as the space of Hochschild $k$-chains
of $\CI(X)$ the space $C_k=\CI(X^{k+1})$, a completion of
$\CI(X)^{\otimes(k+1)}$ (see \cref{hochschild:sec:topological}), with the
face maps \eqref{hochschild:eq:faces-fn}; these are the continuous extensions
of \eqref{hochschild:eq:faces}, and $b=\sum_i(-1)^id_i$ as before. The map
$\varepsilon:C_k\to\Omega^k(X)=\CI(X;\Lambda^kT^*X)$ is defined by
\eqref{hochschild:eq:HKR} in local coordinates. This is invariant: the
mixed derivative $\pa^k u/\pa z_1^{j_1}\cdots\pa z_k^{j_k}$, in which each
variable is differentiated once, restricted to the diagonal, transforms as a
section of $(T^*X)^{\otimes k}$ under change of coordinates, and
$\varepsilon(u)$ is $\frac1{k!}$ times its image under
$v_1\otimes\dots\otimes v_k\mapsto v_1\wedge\dots\wedge v_k$.

\begin{lemma}\label{hochschild:lem:HKR-smooth}
$\varepsilon\circ b=0$ on $C_{k+1}$, and $\varepsilon_*:\HH_k(\CI(X))\to
\Omega^k(X)$ is surjective.
\end{lemma}
\begin{proof}
The value of $\varepsilon(bu)$ at $x\in X$ only involves derivatives of $u$
of order at most $k+1$ at $(x,\dots,x)$. Replacing $u$, in local coordinates,
by its Taylor polynomial of that order, which is a sum of products of
functions of the individual variables, the first statement follows from
\cref{hochschild:lem:HKR-chain}. Embedding $X$ in some $\R^N$, every
$k$-form is a finite sum of forms $a_0\,da_1\wedge\dots\wedge da_k$ with
$a_j\in\CI(X)$ (take $a_1,\dots,a_k$ among the restricted coordinate
functions). By \cref{hochschild:lem:HKR-chain},
$\varphi(a_0;a_1,\dots,a_k)$ is a cycle with $\varepsilon$-image
$a_0\,da_1\wedge\dots\wedge da_k$.
\end{proof}

The main result of this section is the corresponding injectivity.

\begin{theorem}[Connes]\label{hochschild:thm:Connes}
For a compact manifold $X$, $\varepsilon$ induces an isomorphism
$\varepsilon_*:\HH_k(\CI(X))\longrightarrow\Omega^k(X)$ for each $k\in\N_0$.
\end{theorem}

The bar complex $(C_\bullet,b')$ is acyclic by the continuous version of
\cref{hochschild:lem:bar-acyclic}, with $(su)(z_0,\dots,z_{k+1})=u(z_1,\dots,z_{k+1})$;
so $\CI(X)$ is (continuously) H-unital.

\subsection{Reduction to normalized chains}
The following discussion applies verbatim to $\CI(X)$ and to the polynomial
algebra (with $X=\R^n$ and $C_k$ the polynomials on $(\R^n)^{k+1}$). For
$u\in C_k$ consider the diagonals
\begin{equation*}
  D_i=\{z\in X^{k+1};z_i=z_{i+1}\},\ 0\le i\le k-1,\qquad
  D_k=\{z\in X^{k+1};z_k=z_0\},
\end{equation*}
so that $d_iu$ is essentially the restriction of $u$ to $D_i$, and set
\begin{equation*}
\begin{gathered}
  N_k=\{u\in C_k;\ u|_{D_i}=0,\ 0\le i\le k-1\},\qquad
  M_k=\{u\in N_k;\ u|_{D_k}=0\},\\
  B_k=d_{k+1}(N_{k+1})=\{v(z_0,\dots,z_k,z_0);\ v\in N_{k+1}\}.
\end{gathered}
\end{equation*}
Note $N_0=C_0$, $M_0=C_0$, and $B_k\subset M_k$: a $v\in N_{k+1}$ vanishes on
$\{z_k=z_{k+1}\}$, so $v(z_0,\dots,z_k,z_0)$ vanishes when $z_k=z_0$.
If $u\in N_k$ then $d_iu=0$ for $i<k$, so $bu=(-1)^kd_ku$; and $d_ku$ vanishes
on the diagonals $\{z_i=z_{i+1}\}$, $i\le k-2$, of $X^k$. Thus $(N_\bullet,b)$
is a subcomplex of $(C_\bullet,b)$, and
\begin{equation}
  H_k(N_\bullet,b)=M_k/B_k .
  \label{hochschild:eq:HN}
\end{equation}

\begin{proposition}\label{hochschild:prop:normalized}
The inclusion $N_\bullet\hookrightarrow C_\bullet$ induces an isomorphism in
homology. Hence $\HH_k=M_k/B_k$; in particular a chain in $M_k$ which is a
$b$-boundary in $C_\bullet$ lies in $B_k$.
\end{proposition}
\begin{proof}
For $j\in\N_0$ let $J^j_m=\{u\in C_m;\ u|_{D_i}=0\text{ for }0\le i<\min(j,m)\}$.
Thus $J^0=C$, and $J^j_m=N_m$ for $m\le j$. Each $J^j_\bullet$ is a
subcomplex: if $u\in J^j_m$ then $d_iu=0$ for $i<\min(j,m)$, and for the other
$i$ the restriction $d_iu$ still vanishes on $\{z_l=z_{l+1}\}$ for
$l<\min(j,m-1)$ (for $i\le m-1$ because $i\ge l+1$; for $i=m$ trivially).

Let $V^j_m=\{u\in J^j_m;\ u\text{ independent of }z_{j+1}\}$ for $m\ge j+1$,
and $V^j_m=0$ for $m\le j$. For $m\ge j+1$ every $u\in J^j_m$ decomposes
uniquely as
\begin{equation*}
  u=u'+u'',\qquad u'(z)=u(z_0,\dots,z_j,z_j,z_{j+2},\dots,z_m)\in V^j_m,\quad
  u''=u-u'\in J^{j+1}_m,
\end{equation*}
($u''$ vanishes on $\{z_j=z_{j+1}\}$; uniqueness because a function
independent of $z_{j+1}$ vanishing at $z_{j+1}=z_j$ vanishes). So
$J^j_\bullet=J^{j+1}_\bullet\oplus V^j_\bullet$ as graded spaces. We show that
$V^j_\bullet$ is a contractible subcomplex; then this is a direct sum of
complexes and $J^{j+1}\hookrightarrow J^j$ is a quasi-isomorphism.

Let $u\in V^j_m$. In $bu$ the terms $d_iu$, $i<j$, vanish, and
$d_ju=d_{j+1}u$ (both delete the variable $z_{j+1}$ on which $u$ does not
depend; if $j+1=m$, $d_{j+1}$ is the cyclic face $d_m$ and the same holds), so
these cancel. Hence
\begin{equation*}
  bu=\sum_{i=j+2}^m(-1)^id_iu,
\end{equation*}
which is independent of $z_{j+1}$ and lies in $J^j_{m-1}$; so $bu\in
V^j_{m-1}$ (and $bu=0$ if $m=j+1$). Define $H:V^j_m\to V^j_{m+1}$ by inserting
a dummy variable in position $j+2$:
\begin{equation*}
  (Hu)(z_0,\dots,z_{m+1})=u(z_0,\dots,z_{j+1},z_{j+3},\dots,z_{m+1}).
\end{equation*}
Then $d_{j+2}Hu=u$ (if $m=j+1$ this is the cyclic face, which also just
deletes the dummy variable), $d_iHu=Hd_{i-1}u$ for $j+3\le i\le m+1$ (for
$i=m+1$ both sides equal $u(z_0,\dots,z_{j+1},z_{j+3},\dots,z_m,z_0)$ after
renaming), hence
\begin{equation*}
  bHu=(-1)^ju+\sum_{i=j+3}^{m+1}(-1)^iHd_{i-1}u=(-1)^ju-Hbu,
\end{equation*}
i.e.\ $bH+Hb=(-1)^j\Id$ on $V^j_\bullet$. Thus $V^j_\bullet$ is contractible.

Consequently $H_k(C)=H_k(J^0)\cong H_k(J^1)\cong\dots\cong H_k(J^{k+1})$, and
$J^{k+1}$ coincides with $N$ in degrees $\le k+1$, so $H_k(J^{k+1})=H_k(N)$.
All isomorphisms are induced by inclusions. The last statement follows from
\eqref{hochschild:eq:HN}.
\end{proof}

\subsection{Structure of the normalized cycles}
\begin{lemma}\label{hochschild:lem:B-module}
\begin{enumerate}[label=\textup{(\roman*)}]
\item $B_k$ is a $\CI(X^{k+1})$-submodule of $M_k$.
\item If $f\in\CI(X)$ and $g\in M_k$, then $(f(z_j)-f(z_{j+1}))g\in B_k$ for
$0\le j\le k-1$, and $(f(z_k)-f(z_0))g\in B_k$.
\end{enumerate}
\end{lemma}
\begin{proof}
(i) If $\psi\in\CI(X^{k+1})$ and $v\in N_{k+1}$, then
$\tilde\psi v\in N_{k+1}$, where $\tilde\psi(z_0,\dots,z_{k+1})=\psi(z_0,\dots,z_k)$,
and $d_{k+1}(\tilde\psi v)=\psi\,d_{k+1}v$.

(ii) For $0\le j\le k-1$ set
\begin{equation*}
  w(z_0,\dots,z_{k+1})=(-1)^{j+1}\bigl(f(z_j)-f(z_{j+1})\bigr)\,
  g(z_0,\dots,z_j,z_{j+2},\dots,z_{k+1})\in C_{k+1}.
\end{equation*}
We compute $bw$. For $i<j$, $d_iw=0$ since $g$ vanishes on $\{z_i=z_{i+1}\}$;
$d_jw=0$ because of the first factor; for $j+2\le i\le k$, $d_iw=0$ because
$g$ vanishes on its diagonal $D_{i-1}$; and $d_{k+1}w=0$ because $g$ vanishes
on its diagonal $D_k$ (here $j\le k-1$ is used, so that $z_{k+1}$ is a variable
of $g$). Finally $d_{j+1}w$ identifies $z_{j+1}$ with $z_{j+2}$ and gives
$(-1)^{j+1}(f(z_j)-f(z_{j+1}))g(z_0,\dots,z_k)$. Hence
$bw=(f(z_j)-f(z_{j+1}))g$. This lies in $M_k$ and is a boundary in $C_\bullet$,
so it lies in $B_k$ by \cref{hochschild:prop:normalized}. The last case
follows from $f(z_k)-f(z_0)=-\sum_{j=0}^{k-1}(f(z_j)-f(z_{j+1}))$.
\end{proof}

Near the full diagonal $\Delta=\{z_0=z_1=\dots=z_k\}$ we use the following
coordinates. Let $x_0\in X$, choose coordinates near $x_0$ identifying a
neighbourhood with an open set in $\R^n$, and on a neighbourhood
$W=\{|x-x_0|<\delta,\ |y_j|<\delta\}$ of $(x_0,\dots,x_0)$ use
\begin{equation}
  x=z_0,\qquad y_j=z_j-z_{j-1},\ j=1,\dots,k .
  \label{hochschild:eq:diag-coords}
\end{equation}
Then $D_{j-1}=\{y_j=0\}$ for $1\le j\le k$, $D_k=\{y_1+\dots+y_k=0\}$ and
$\Delta=\{y=0\}$.

\begin{lemma}\label{hochschild:lem:antisym}
Let $q$ be a $k$-linear form on $\R^n$ (with complex values). Then
$q(y_1,\dots,y_{k-1},-(y_1+\dots+y_{k-1}))\equiv0$ if and only if $q$ is
alternating.
\end{lemma}
\begin{proof}
If $q$ is alternating the expression is
$-\sum_{j<k}q(y_1,\dots,y_{k-1},y_j)=0$. Conversely each term
$q(y_1,\dots,y_{k-1},y_j)$, $j<k$, is of degree $2$ in $y_j$ and of degree $1$
in the other $y_i$, $i<k$; these multidegrees are distinct, so each term
vanishes separately. Thus $q$ vanishes whenever the $j$-th and $k$-th
arguments coincide, i.e.\ $q$ is alternating in these two arguments. The
transpositions $(j\ k)$ generate $\mathfrak S_k$, so $q$ is alternating.
\end{proof}

\begin{lemma}\label{hochschild:lem:leading}
Let $u\in N_k$. In the coordinates \eqref{hochschild:eq:diag-coords},
$u=\sum_Iy_{1,i_1}\cdots y_{k,i_k}w_I(x,y)$ on $W$, with $w_I$ smooth, where
$I=(i_1,\dots,i_k)\in\{1,\dots,n\}^k$ and $y_{j,i}$ is the $i$-th component of
$y_j$. Hence the Taylor series of $u$ at $y=0$ (with coefficients smooth in
$x$) starts with the $k$-linear term $u_{(k)}(x;y)=\sum_Iw_I(x,0)y_{1,i_1}\cdots
y_{k,i_k}$, and
\begin{equation*}
  \varepsilon(u)(x)=\frac1{k!}\sum_Iw_I(x,0)\,dx^{i_1}\wedge\dots\wedge dx^{i_k}.
\end{equation*}
If $u\in M_k$, then $u_{(k)}(x;\cdot)$ is alternating and $\varepsilon(u)(x)=0$
if and only if $u_{(k)}(x;\cdot)=0$, i.e.\ if and only if $u$ vanishes to
order $k+1$ at $(x,\dots,x)$.
\end{lemma}
\begin{proof}
By Hadamard's lemma in $y_1$, $u=\sum_iy_{1,i}u_i$ with
$u_i=\int_0^1(\pa_{y_{1,i}}u)(x,ty_1,y_2,\dots,y_k)\,dt$; the $u_i$ still
vanish on $\{y_j=0\}$, $j\ge2$. Repeating in $y_2,\dots,y_k$ gives the
representation. Since $z_l=x+y_1+\dots+y_l$, the derivative $\pa_{z_l}$
equals $\pa_{y_l}-\pa_{y_{l+1}}$ ($\pa_{y_l}$ for $l=k$). In
$\pa_{z_1^{i_1}}\cdots\pa_{z_k^{i_k}}u$ at $y=0$ only terms in which every
factor $y_{l,\cdot}$ is differentiated survive; since $\pa_{z_k}$ can only hit
$y_k$, then $\pa_{z_{k-1}}$ only $y_{k-1}$, and so on, the result is
$w_I(x,0)$, which gives the formula for $\varepsilon(u)$. If $u\in M_k$, the
homogeneous term $u_{(k)}$ vanishes on $\{y_1+\dots+y_k=0\}$, so it is
alternating by \cref{hochschild:lem:antisym}. An alternating tensor is
determined by its image in $\Lambda^k$, whence the last statement.
\end{proof}

\begin{lemma}\label{hochschild:lem:poly-lift}
Let $k\ge1$ and $m\ge k+1$. Let $\mathcal M_m$ be the space of homogeneous
polynomials $q$ of degree $m$ in $y=(y_1,\dots,y_k)\in(\R^n)^k$ which vanish
on each of the subspaces $\{y_j=0\}$, $1\le j\le k$, and
$\{y_1+\dots+y_k=0\}$, and let $\mathcal V_m$ be the space of homogeneous
polynomials of degree $m$ in $(y_1,\dots,y_{k+1})\in(\R^n)^{k+1}$ vanishing on
each $\{y_j=0\}$, $1\le j\le k+1$. Then
\begin{equation*}
  R:\mathcal V_m\longrightarrow\mathcal M_m,\quad
  (RV)(y_1,\dots,y_k)=V(y_1,\dots,y_k,-(y_1+\dots+y_k)),
\end{equation*}
is surjective.
\end{lemma}
\begin{proof}
Given $q\in\mathcal M_m$, the polynomial chain
$u(z_0,\dots,z_k)=q(z_1-z_0,\dots,z_k-z_{k-1})$ for $\C[x]$ lies in $M_k$, and
$\varepsilon(u)=0$ since $u$ vanishes to order $m>k$ on $\Delta$. By
\cref{hochschild:thm:polynomial} and \cref{hochschild:prop:normalized} (for
polynomials) $u\in B_k$, i.e.\ $u=d_{k+1}v$ with $v\in N_{k+1}$. Writing $v$
in the coordinates $x=z_0$, $y_j=z_j-z_{j-1}$, $1\le j\le k+1$, as
$V(x;y_1,\dots,y_{k+1})$, $V$ vanishes on the subspaces $\{y_j=0\}$, and
$d_{k+1}v=V(x;y_1,\dots,y_k,-(y_1+\dots+y_k))$ because $z_{k+1}=z_0$ means
$y_{k+1}=-(y_1+\dots+y_k)$. The component of $V$ of degree $0$ in $x$ and
degree $m$ in $y$ lies in $\mathcal V_m$ and is mapped by $R$ to the
corresponding component of $u$, which is $q$.
\end{proof}

\begin{lemma}\label{hochschild:lem:local}
Let $u\in M_k$, $k\ge1$, satisfy $\varepsilon(u)=0$. Then every point $p\in
X^{k+1}$ has a neighbourhood on which $u$ coincides with an element of $B_k$.
\end{lemma}
\begin{proof}
First let $p\notin\Delta$. Then $p_j\ne p_{j+1}$ for some $0\le j\le k-1$ (else
all $p_j$ coincide), and we choose $f\in\CI(X)$ with $f(p_j)\ne f(p_{j+1})$
and $\psi\in\CI(X^{k+1})$ supported where $f(z_j)\ne f(z_{j+1})$, equal to
$1$ near $p$. Then $g=\psi u/(f(z_j)-f(z_{j+1}))\in M_k$ and $u=(f(z_j)-f(z_{j+1}))g$
near $p$; this lies in $B_k$ by \cref{hochschild:lem:B-module}.

Now let $p=(x_0,\dots,x_0)\in\Delta$ and use the coordinates
\eqref{hochschild:eq:diag-coords}. The Taylor series of $u$ in $y$ at $y=0$ is
$\sum_mu_{(m)}(x;y)$, where $u_{(m)}(x;\cdot)$ is homogeneous of degree $m$
and vanishes on the same linear subspaces as $u$, so
$u_{(m)}\in\CI(B)\otimes\mathcal M_m$, $B=\{|x-x_0|<\delta\}$. By
\cref{hochschild:lem:leading}, $u_{(m)}=0$ for $m\le k$. For $m\ge k+1$,
choosing a linear right inverse of $R$ in \cref{hochschild:lem:poly-lift}, we
find $V_{(m)}\in\CI(B)\otimes\mathcal V_m$ with $RV_{(m)}=u_{(m)}$. By Borel's
lemma (see \cite[Section~1.2]{Hormander-I}), shrinking $\delta$, there is a smooth
$V(x;y_1,\dots,y_{k+1})=\sum_m\chi(\lambda_m|y|)V_{(m)}(x;y)$, with
$\chi\in\CIc(\R)$ equal to $1$ near $0$ and $\lambda_m\to\infty$ fast
enough, having Taylor series $\sum_mV_{(m)}$ at $y=0$; each term, hence $V$,
vanishes on the subspaces $\{y_j=0\}$. Let $\Theta\in\CIc$ be a cut-off near
$(x_0,\dots,x_0)\in X^{k+2}$ and let $v=\Theta V\in N_{k+1}$, extended by
zero (written in the variables $x=z_0$, $y_j=z_j-z_{j-1}$, $j\le k+1$). Then
$s=d_{k+1}v\in B_k$ has the same Taylor series as $u$ at $\Delta$ near $p$. So
$g=u-s\in M_k$ vanishes to infinite order at $\Delta$ near $p$. Near $p$,
\begin{equation*}
  g=\sum_{j=1}^k\sum_{i=1}^ny_{j,i}\,g_{j,i},\qquad
  g_{j,i}=\frac{y_{j,i}\,g}{|y|^2},
\end{equation*}
where $g_{j,i}$, extended by $0$ on $\Delta$, is smooth since $g$ is flat
there, and vanishes wherever $g$ does. Taking $f_i\in\CI(X)$ equal to the
$i$-th coordinate near $x_0$, and a cut-off $\psi$ equal to $1$ near $p$, we
have $y_{j,i}=-(f_i(z_{j-1})-f_i(z_j))$ on the support of $\psi$, so near $p$,
$g=-\sum_{j,i}(f_i(z_{j-1})-f_i(z_j))\psi g_{j,i}\in B_k$ by
\cref{hochschild:lem:B-module}. Hence $u=s+g$ coincides near $p$ with an
element of $B_k$.
\end{proof}

\begin{proof}[Proof of \cref{hochschild:thm:Connes}]
For $k=0$, $(bu)(z_0)=u(z_0,z_0)-u(z_0,z_0)=0$ for $u\in C_1$, so
$\HH_0=\CI(X)=\Omega^0(X)$. Let $k\ge1$. By
\cref{hochschild:lem:HKR-smooth} it remains to show injectivity. By
\cref{hochschild:prop:normalized} every class in $\HH_k$ is represented by
some $u\in M_k$, and $u$ is a boundary if (and only if) $u\in B_k$. Suppose
$\varepsilon(u)=0$. By \cref{hochschild:lem:local} and compactness there is a
finite partition of unity $\sum_\beta\psi_\beta=1$ on $X^{k+1}$ and
$s_\beta\in B_k$ with $u=s_\beta$ on $\supp\psi_\beta$. Then
$u=\sum_\beta\psi_\beta s_\beta\in B_k$ by \cref{hochschild:lem:B-module}(i).
\end{proof}

\begin{remark}\label{hochschild:rem:noncompact}
Compactness was only used to make the partitions of unity finite. With some
care about supports (using locally finite partitions of unity and cut-offs
in all variables) the same argument applies to non-compact $X$.
\end{remark}
